% Options for packages loaded elsewhere
\PassOptionsToPackage{unicode}{hyperref}
\PassOptionsToPackage{hyphens}{url}
\PassOptionsToPackage{dvipsnames,svgnames,x11names}{xcolor}
\documentclass[
  11pt,
]{article}
\usepackage{xcolor}
\usepackage[margin=2cm]{geometry}
\usepackage{amsmath,amssymb}
\setcounter{secnumdepth}{-\maxdimen} % remove section numbering
\usepackage{iftex}
\ifPDFTeX
  \usepackage[T1]{fontenc}
  \usepackage[utf8]{inputenc}
  \usepackage{textcomp} % provide euro and other symbols
\else % if luatex or xetex
  \usepackage{unicode-math} % this also loads fontspec
  \defaultfontfeatures{Scale=MatchLowercase}
  \defaultfontfeatures[\rmfamily]{Ligatures=TeX,Scale=1}
\fi
\usepackage{lmodern}
\ifPDFTeX\else
  % xetex/luatex font selection
    \setmonofont[Scale=0.8,BoldFont=DejaVuSansMono-Bold.ttf]{DejaVuSansMono.ttf}
\fi
% Use upquote if available, for straight quotes in verbatim environments
\IfFileExists{upquote.sty}{\usepackage{upquote}}{}
\IfFileExists{microtype.sty}{% use microtype if available
  \usepackage[]{microtype}
  \UseMicrotypeSet[protrusion]{basicmath} % disable protrusion for tt fonts
}{}
\makeatletter
\@ifundefined{KOMAClassName}{% if non-KOMA class
  \IfFileExists{parskip.sty}{%
    \usepackage{parskip}
  }{% else
    \setlength{\parindent}{0pt}
    \setlength{\parskip}{6pt plus 2pt minus 1pt}}
}{% if KOMA class
  \KOMAoptions{parskip=half}}
\makeatother
\usepackage{color}
\usepackage{fancyvrb}
\newcommand{\VerbBar}{|}
\newcommand{\VERB}{\Verb[commandchars=\\\{\}]}
\DefineVerbatimEnvironment{Highlighting}{Verbatim}{commandchars=\\\{\}}
% Add ',fontsize=\small' for more characters per line
\newenvironment{Shaded}{}{}
\newcommand{\AlertTok}[1]{\textcolor[rgb]{1.00,0.00,0.00}{\textbf{#1}}}
\newcommand{\AnnotationTok}[1]{\textcolor[rgb]{0.38,0.63,0.69}{\textbf{\textit{#1}}}}
\newcommand{\AttributeTok}[1]{\textcolor[rgb]{0.49,0.56,0.16}{#1}}
\newcommand{\BaseNTok}[1]{\textcolor[rgb]{0.25,0.63,0.44}{#1}}
\newcommand{\BuiltInTok}[1]{\textcolor[rgb]{0.00,0.50,0.00}{#1}}
\newcommand{\CharTok}[1]{\textcolor[rgb]{0.25,0.44,0.63}{#1}}
\newcommand{\CommentTok}[1]{\textcolor[rgb]{0.38,0.63,0.69}{\textit{#1}}}
\newcommand{\CommentVarTok}[1]{\textcolor[rgb]{0.38,0.63,0.69}{\textbf{\textit{#1}}}}
\newcommand{\ConstantTok}[1]{\textcolor[rgb]{0.53,0.00,0.00}{#1}}
\newcommand{\ControlFlowTok}[1]{\textcolor[rgb]{0.00,0.44,0.13}{\textbf{#1}}}
\newcommand{\DataTypeTok}[1]{\textcolor[rgb]{0.56,0.13,0.00}{#1}}
\newcommand{\DecValTok}[1]{\textcolor[rgb]{0.25,0.63,0.44}{#1}}
\newcommand{\DocumentationTok}[1]{\textcolor[rgb]{0.73,0.13,0.13}{\textit{#1}}}
\newcommand{\ErrorTok}[1]{\textcolor[rgb]{1.00,0.00,0.00}{\textbf{#1}}}
\newcommand{\ExtensionTok}[1]{#1}
\newcommand{\FloatTok}[1]{\textcolor[rgb]{0.25,0.63,0.44}{#1}}
\newcommand{\FunctionTok}[1]{\textcolor[rgb]{0.02,0.16,0.49}{#1}}
\newcommand{\ImportTok}[1]{\textcolor[rgb]{0.00,0.50,0.00}{\textbf{#1}}}
\newcommand{\InformationTok}[1]{\textcolor[rgb]{0.38,0.63,0.69}{\textbf{\textit{#1}}}}
\newcommand{\KeywordTok}[1]{\textcolor[rgb]{0.00,0.44,0.13}{\textbf{#1}}}
\newcommand{\NormalTok}[1]{#1}
\newcommand{\OperatorTok}[1]{\textcolor[rgb]{0.40,0.40,0.40}{#1}}
\newcommand{\OtherTok}[1]{\textcolor[rgb]{0.00,0.44,0.13}{#1}}
\newcommand{\PreprocessorTok}[1]{\textcolor[rgb]{0.74,0.48,0.00}{#1}}
\newcommand{\RegionMarkerTok}[1]{#1}
\newcommand{\SpecialCharTok}[1]{\textcolor[rgb]{0.25,0.44,0.63}{#1}}
\newcommand{\SpecialStringTok}[1]{\textcolor[rgb]{0.73,0.40,0.53}{#1}}
\newcommand{\StringTok}[1]{\textcolor[rgb]{0.25,0.44,0.63}{#1}}
\newcommand{\VariableTok}[1]{\textcolor[rgb]{0.10,0.09,0.49}{#1}}
\newcommand{\VerbatimStringTok}[1]{\textcolor[rgb]{0.25,0.44,0.63}{#1}}
\newcommand{\WarningTok}[1]{\textcolor[rgb]{0.38,0.63,0.69}{\textbf{\textit{#1}}}}
\usepackage{longtable,booktabs,array}
\usepackage{calc} % for calculating minipage widths
% Correct order of tables after \paragraph or \subparagraph
\usepackage{etoolbox}
\makeatletter
\patchcmd\longtable{\par}{\if@noskipsec\mbox{}\fi\par}{}{}
\makeatother
% Allow footnotes in longtable head/foot
\IfFileExists{footnotehyper.sty}{\usepackage{footnotehyper}}{\usepackage{footnote}}
\makesavenoteenv{longtable}
\ifLuaTeX
\usepackage[bidi=basic]{babel}
\else
\usepackage[bidi=default]{babel}
\fi
\babelprovide[main,import]{french}
% get rid of language-specific shorthands (see #6817):
\let\LanguageShortHands\languageshorthands
\def\languageshorthands#1{}
\setlength{\emergencystretch}{3em} % prevent overfull lines
\providecommand{\tightlist}{%
  \setlength{\itemsep}{0pt}\setlength{\parskip}{0pt}}
% Ajouté au préambule de chaque document par generer.sh
\usepackage{xcolor}
\usepackage{newunicodechar}
\newfontfamily\symbolfont{DejaVuSans.ttf}
% Caractères absents de Latin Modern (flèches, comparateurs, cases) : repli sur DejaVu Sans
\newunicodechar{→}{{\symbolfont →}}
\newunicodechar{←}{{\symbolfont ←}}
\newunicodechar{⇄}{{\symbolfont ⇄}}
\newunicodechar{↔}{{\symbolfont ↔}}
\newunicodechar{≥}{{\symbolfont ≥}}
\newunicodechar{≤}{{\symbolfont ≤}}
\newunicodechar{≠}{{\symbolfont ≠}}
\newunicodechar{≈}{{\symbolfont ≈}}
\newunicodechar{✓}{{\symbolfont ✓}}
\newunicodechar{✗}{{\symbolfont ✗}}
\newunicodechar{⚠}{{\symbolfont ⚠}}
\newunicodechar{✅}{{\symbolfont\color[rgb]{0.1,0.55,0.2}✔}}
\newunicodechar{‑}{\mbox{-}}
% Retour à la ligne dans les blocs de code et les URL longues
\usepackage{fvextra}
\DefineVerbatimEnvironment{Highlighting}{Verbatim}{breaklines,commandchars=\\\{\}}
\usepackage{xurl}
\setlength{\emergencystretch}{3em}
\AtBeginEnvironment{longtable}{\small}
% Figures flottantes (pas de demi-pages vides), jamais plus hautes que 85 % de la page
\usepackage{float}
\floatplacement{figure}{htbp}
\usepackage{adjustbox}
\providecommand*\pandocbounded[1]{#1}% absent des documents sans image
\renewcommand*\pandocbounded[1]{\adjustbox{max width=\linewidth,max totalheight=0.85\textheight}{#1}}
\usepackage[font=small,labelfont=bf]{caption}
% Contenu sur une nouvelle page après la table des matières
\let\tdmorigine\tableofcontents
\renewcommand{\tableofcontents}{\tdmorigine\newpage}
\usepackage{fancyhdr}
\pagestyle{fancy}
\fancyhf{}
\fancyhead[L]{\small SmartRecruit — Qualité logicielle ISO/IEC 25010}
\fancyhead[R]{\small\thepage}
\renewcommand{\headrulewidth}{0.4pt}
\usepackage{bookmark}
\IfFileExists{xurl.sty}{\usepackage{xurl}}{} % add URL line breaks if available
\urlstyle{same}
\hypersetup{
  pdftitle={Fiche d\textquotesingle étude --- Membre 3 : performance, fiabilité, portabilité},
  pdflang={fr},
  colorlinks=true,
  linkcolor={Maroon},
  filecolor={Maroon},
  citecolor={Blue},
  urlcolor={Blue},
  pdfcreator={LaTeX via pandoc}}

\title{Fiche d\textquotesingle étude --- Membre 3 : performance,
fiabilité, portabilité}
\author{}
\date{}

\begin{document}
\maketitle

\renewcommand*\contentsname{Table des matières}
{
\hypersetup{linkcolor=}
\setcounter{tocdepth}{3}
\tableofcontents
}
\begin{quote}
Pour : \emph{(à compléter)}. Complète le \href{SOUTENANCE.pdf}{guide de
soutenance} (questions-réponses, membre 3) et la
\href{REPARTITION_TESTS.pdf}{répartition des tests}.
\end{quote}

\section{1. Ta mission en soutenance}\label{1-ta-mission-en-soutenance}

Le lead lance k6 en direct ; toi, tu dois pouvoir \textbf{lire le
résultat à voix haute}, \textbf{relancer} tes tests et \textbf{répondre}
sur :

\begin{itemize}
\tightlist
\item
  \textbf{6.2 Efficacité de performance} (temps de réponse sous charge)
  ;
\item
  \textbf{6.5 Fiabilité} (comportement en cas de panne, pas de perte de
  données) ;
\item
  \textbf{6.8 Portabilité} (installation reproductible, CI/CD) ;
\item
  la compatibilité avec un système tiers (API LLM externe).
\end{itemize}

\section{2. Ce que tu dois apprendre
(notions)}\label{2-ce-que-tu-dois-apprendre-notions}

\begin{longtable}[]{@{}
  >{\raggedright\arraybackslash}p{(\linewidth - 2\tabcolsep) * \real{0.4309}}
  >{\raggedright\arraybackslash}p{(\linewidth - 2\tabcolsep) * \real{0.5691}}@{}}
\toprule\noalign{}
\begin{minipage}[b]{\linewidth}\raggedright
Notion
\end{minipage} & \begin{minipage}[b]{\linewidth}\raggedright
Ce que tu dois savoir dire
\end{minipage} \\
\midrule\noalign{}
\endhead
\bottomrule\noalign{}
\endlastfoot
Test de charge / stress / endurance & Charge : comportement sous un
trafic attendu (notre cas, 20 utilisateurs). Stress : on augmente
jusqu\textquotesingle à la rupture. Endurance : charge longue durée
(fuites mémoire). \\
Utilisateur virtuel (VU) & Un script qui rejoue un utilisateur en boucle
; k6 en fait tourner 20 en parallèle. \\
Débit & Requêtes par seconde (req/s). \\
Latence, médiane, percentile & p95 = 95 \% des requêtes sont plus
rapides que cette valeur. La moyenne cache les requêtes lentes ; le p95
les fait apparaître. \\
Seuil (threshold) & Critère d\textquotesingle acceptation codé dans le
script (\texttt{p(95)\textless{}500}) : k6 échoue s\textquotesingle il
n\textquotesingle est pas tenu, donc le verdict est automatique. \\
Fiabilité, tolérance aux pannes & Le système continue de rendre service
quand un composant tombe, et revient à l\textquotesingle état normal
ensuite. \\
« Best effort » & Un appel non indispensable (l\textquotesingle IA) peut
échouer sans bloquer l\textquotesingle action principale (la candidature
reste enregistrée, score \texttt{null}). \\
Délai d\textquotesingle attente (timeout) & Ne pas attendre indéfiniment
un service arrêté : le gateway abandonne après 3 s et répond une erreur
claire. \\
Reprise automatique (retry) & Réessayer plus tard un traitement échoué
(relance de l\textquotesingle analyse IA des offres chaque minute). \\
Test de chaos ciblé & Provoquer volontairement une vraie panne (arrêt
d\textquotesingle un conteneur) et vérifier le comportement. \\
Durabilité des données & Une transaction validée survit au redémarrage
de la base (PostgreSQL + volume Docker). \\
Conteneur, image, Dockerfile, Compose & Image = application +
dépendances figées ; conteneur = image en exécution ; Compose démarre
tous les conteneurs ensemble. \\
CI/CD & À chaque push, GitHub Actions reconstruit et teste tout (CI),
puis déploie (CD) sur la branche \texttt{main}. \\
Kubernetes, kind, Helm, Kustomize, Terraform, Ansible & Kubernetes
orchestre les conteneurs (kind = Kubernetes dans Docker) ; Helm et
Kustomize décrivent le déploiement ; Terraform crée
l\textquotesingle infrastructure ; Ansible configure la VM. \\
Test de fumée (smoke test) & Vérification rapide après déploiement : les
pages et l\textquotesingle API principales répondent. \\
\end{longtable}

\section{3. Fichiers à lire, dans
l\textquotesingle ordre}\label{3-fichiers-uxe0-lire-dans-lordre}

\begin{longtable}[]{@{}
  >{\raggedright\arraybackslash}p{(\linewidth - 4\tabcolsep) * \real{0.0155}}
  >{\raggedright\arraybackslash}p{(\linewidth - 4\tabcolsep) * \real{0.7674}}
  >{\raggedright\arraybackslash}p{(\linewidth - 4\tabcolsep) * \real{0.2171}}@{}}
\toprule\noalign{}
\begin{minipage}[b]{\linewidth}\raggedright
\#
\end{minipage} & \begin{minipage}[b]{\linewidth}\raggedright
Fichier
\end{minipage} & \begin{minipage}[b]{\linewidth}\raggedright
Ce que tu y cherches
\end{minipage} \\
\midrule\noalign{}
\endhead
\bottomrule\noalign{}
\endlastfoot
1 &
\texttt{quality/}\allowbreak{}\texttt{perf/}\allowbreak{}\texttt{load-}\allowbreak{}\texttt{test.}\allowbreak{}\texttt{js}
& Le scénario \texttt{constant-vus} (20 VU, 1 min), le mélange des 3
rôles, le seuil \texttt{p(95)\textless{}500}, le résumé imprimé. \\
2 &
\texttt{quality/}\allowbreak{}\texttt{reports/}\allowbreak{}\texttt{perf/}\allowbreak{}\texttt{k6-}\allowbreak{}\texttt{console.}\allowbreak{}\texttt{txt}
& La sortie réelle : requêtes, débit, médiane, p95, max, taux
d\textquotesingle erreur, seuils \texttt{OK}. \\
3 &
\texttt{quality/}\allowbreak{}\texttt{reliability/}\allowbreak{}\texttt{test\_}\allowbreak{}\texttt{reliability.}\allowbreak{}\texttt{py}
& Les \textbf{8 scénarios} (noms en français, lisibles) et comment ils
arrêtent/redémarrent les conteneurs. \\
4 &
\texttt{backend/}\allowbreak{}\texttt{job-}\allowbreak{}\texttt{service/}\allowbreak{}\texttt{src/}\allowbreak{}\texttt{main/}\allowbreak{}\texttt{java/}\allowbreak{}\texttt{com/}\allowbreak{}\texttt{smartrecruit/}\allowbreak{}\texttt{job/}\allowbreak{}\texttt{service/}\allowbreak{}\texttt{Job}\allowbreak{}\texttt{Analysis}\allowbreak{}\texttt{Listener.}\allowbreak{}\texttt{java}
& Correctif du défaut n°1 : relance automatique des analyses IA
échouées. \\
5 &
\texttt{backend/}\allowbreak{}\texttt{api-}\allowbreak{}\texttt{gateway/}\allowbreak{}\texttt{src/}\allowbreak{}\texttt{main/}\allowbreak{}\texttt{java/}\allowbreak{}\texttt{com/}\allowbreak{}\texttt{smartrecruit/}\allowbreak{}\texttt{gateway/}\allowbreak{}\texttt{config/}\allowbreak{}\texttt{Gateway}\allowbreak{}\texttt{Config.}\allowbreak{}\texttt{java}
& Correctif du défaut n°2 : délai de connexion borné à 3 s. \\
6 &
\texttt{backend/}\allowbreak{}\texttt{application-}\allowbreak{}\texttt{service/}\allowbreak{}\texttt{src/}\allowbreak{}\texttt{main/}\allowbreak{}\texttt{java/}\allowbreak{}\texttt{com/}\allowbreak{}\texttt{smartrecruit/}\allowbreak{}\texttt{application/}\allowbreak{}\texttt{client/}\allowbreak{}\texttt{Ai}\allowbreak{}\texttt{Matching}\allowbreak{}\texttt{Client.}\allowbreak{}\texttt{java}
& L\textquotesingle appel à l\textquotesingle IA en « best effort ». \\
7 & \texttt{docker-compose.yml} et le
\href{GUIDE_INSTALLATION.pdf}{guide d\textquotesingle installation} &
Les services, les ports, les volumes, l\textquotesingle installation par
un tiers. \\
8 &
\texttt{.}\allowbreak{}\texttt{github/}\allowbreak{}\texttt{workflows/}\allowbreak{}\texttt{pipeline.}\allowbreak{}\texttt{yml}
& Les jobs : \texttt{backend}, \texttt{frontend}, \texttt{ai-service},
\texttt{kustomize}, \texttt{helm}, \texttt{terraform}, \texttt{ansible},
\texttt{images}, \texttt{deploy}, \texttt{smoke-test}. \\
9 & \texttt{deploy/} (\texttt{k8s/}, \texttt{helm/},
\texttt{terraform/}, \texttt{ansible/}) et \texttt{docs/DEVOPS.md} &
Survol : savoir dire à quoi sert chaque dossier. \\
10 &
\texttt{ai-}\allowbreak{}\texttt{service/}\allowbreak{}\texttt{tests/}\allowbreak{}\texttt{test\_}\allowbreak{}\texttt{llm\_}\allowbreak{}\texttt{client.}\allowbreak{}\texttt{py}
& Coexistence avec un système tiers (API OpenRouter) sans dépendre du
réseau. \\
\end{longtable}

\section{4. Commandes à savoir lancer (et
lire)}\label{4-commandes-uxe0-savoir-lancer-et-lire}

Application démarrée au préalable
(\texttt{docker\ compose\ up\ -}\allowbreak{}\texttt{d\ -}\allowbreak{}\texttt{-}\allowbreak{}\texttt{build}).

\begin{Shaded}
\begin{Highlighting}[]
\FunctionTok{bash}\NormalTok{ quality/run{-}all.sh perf}
\end{Highlighting}
\end{Shaded}

Charge k6 (1 min). Lire : nombre de requêtes, req/s, p95,
\texttt{http\_req\_failed}, puis \texttt{seuil\ …\ :\ OK}.

\begin{Shaded}
\begin{Highlighting}[]
\FunctionTok{bash}\NormalTok{ quality/run{-}all.sh reliability}
\end{Highlighting}
\end{Shaded}

Les 8 pannes réelles (plusieurs minutes : les conteneurs
s\textquotesingle arrêtent et redémarrent). Lire : \texttt{8\ passed}.

\begin{Shaded}
\begin{Highlighting}[]
\FunctionTok{bash}\NormalTok{ quality/run{-}all.sh install}
\end{Highlighting}
\end{Shaded}

Reconstruit et démarre toute la stack depuis le dépôt : la preuve de
portabilité sur le poste.

\section{5. Chiffres à connaître par
cœur}\label{5-chiffres-uxe0-connauxeetre-par-cux153ur}

\begin{longtable}[]{@{}
  >{\raggedright\arraybackslash}p{(\linewidth - 2\tabcolsep) * \real{0.2821}}
  >{\raggedright\arraybackslash}p{(\linewidth - 2\tabcolsep) * \real{0.7179}}@{}}
\toprule\noalign{}
\begin{minipage}[b]{\linewidth}\raggedright
Sujet
\end{minipage} & \begin{minipage}[b]{\linewidth}\raggedright
Valeur
\end{minipage} \\
\midrule\noalign{}
\endhead
\bottomrule\noalign{}
\endlastfoot
Charge & \textbf{20} utilisateurs simultanés pendant \textbf{1 min}, sur
\textbf{3 068} offres \\
Résultat & \textbf{4 463} requêtes (\textbf{72 req/s}), p95
\textbf{349,6 ms} (seuil 500), médiane \textbf{55 ms}, max ≈ 1,5 s,
\textbf{0,00 \%} d\textquotesingle erreurs \\
Fiabilité & \textbf{8} scénarios ; premier passage \textbf{5/8},
\textbf{2 défauts} trouvés et corrigés, final \textbf{8/8} \\
CI & \textbf{20/20} jobs sur la version finale (commit
\texttt{ed0779b}) \\
Déploiement & VM Azure + tests de fumée : \textbf{22/22} jobs (commit
\texttt{3b4cd95}) \\
Images Docker & \textbf{8} (6 services Java, service IA, frontend) \\
\end{longtable}

D\textquotesingle où viennent 20 et 22 jobs : 6 services backend +
frontend + IA + kustomize + helm + terraform + ansible + 8 images = 20 ;
avec \texttt{deploy} et \texttt{smoke-test} = 22.

\section{6. Preuves à savoir ouvrir}\label{6-preuves-uxe0-savoir-ouvrir}

\begin{itemize}
\tightlist
\item
  \texttt{quality/}\allowbreak{}\texttt{reports/}\allowbreak{}\texttt{perf/}\allowbreak{}\texttt{k6-}\allowbreak{}\texttt{console.}\allowbreak{}\texttt{txt}
  et \texttt{k6-summary.json}.
\item
  \texttt{quality/}\allowbreak{}\texttt{reports/}\allowbreak{}\texttt{reliability/}\allowbreak{}\texttt{junit.}\allowbreak{}\texttt{xml}
  et \texttt{pytest-console.txt}.
\item
  Les exécutions GitHub Actions liées dans la
  \href{MATRICE_TRACABILITE.pdf}{matrice} (ligne EQ‑08) et le
  \href{RAPPORT_CONFORMITE_ISO25010.pdf}{rapport de conformité}.
\end{itemize}

\section{7. Questions pièges (au-delà du guide de
soutenance)}\label{7-questions-piuxe8ges-au-deluxe0-du-guide-de-soutenance}

\begin{longtable}[]{@{}
  >{\raggedright\arraybackslash}p{(\linewidth - 2\tabcolsep) * \real{0.5000}}
  >{\raggedright\arraybackslash}p{(\linewidth - 2\tabcolsep) * \real{0.5000}}@{}}
\toprule\noalign{}
\begin{minipage}[b]{\linewidth}\raggedright
Question
\end{minipage} & \begin{minipage}[b]{\linewidth}\raggedright
Réponse
\end{minipage} \\
\midrule\noalign{}
\endhead
\bottomrule\noalign{}
\endlastfoot
Pourquoi 20 utilisateurs et 500 ms ? & C\textquotesingle est
l\textquotesingle exemple chiffré du cahier des charges (§6.2) ; on
l\textquotesingle a pris comme exigence mesurable et codé comme
seuil. \\
Votre test de charge est-il réaliste ? & Il mélange les 3 rôles et les
actions courantes (recherche, détail, candidatures, statistiques, audit)
sur un fort volume. Limite : générateur de charge et application sur la
même machine. \\
Comment le test de charge passe-t-il l\textquotesingle anti-force brute
(10 connexions / 5 min / IP) ? & Le script ne se connecte que \textbf{3
fois}, dans \texttt{setup()} (un compte par rôle) ; les 20 utilisateurs
virtuels réutilisent ces jetons. On mesure donc la charge réelle de
l\textquotesingle application, pas celle de la connexion. \\
Qu\textquotesingle est-ce qui prouve « pas de perte de données » ? & Le
scénario PostgreSQL : création d\textquotesingle une offre,
d\textquotesingle une candidature et d\textquotesingle un statut,
redémarrage de la base, relecture : tout est identique. \\
Et si le service candidatures tombe pendant qu\textquotesingle un
utilisateur postule ? & Le gateway répond vite une erreur (pas
d\textquotesingle attente de 30 s) ; après redémarrage, le service
reprend sans perte (scénario 7). \\
Les 2 défauts trouvés auraient-ils été vus en production ? & Oui : une
offre créée pendant une panne de l\textquotesingle IA
n\textquotesingle était jamais analysée, et un service arrêté bloquait
les utilisateurs plus de 30 s. Les tests les ont révélés avant. \\
La portabilité est-elle prouvée sur une autre machine ? & Oui par la CI
: un runner Ubuntu neuf reconstruit tout à chaque push, et le
déploiement part de zéro sur une VM Azure. Idéalement, faire aussi
installer un membre sur son PC en suivant le guide. \\
Docker ou machine virtuelle ? & Un conteneur partage le noyau de
l\textquotesingle hôte : plus léger et plus rapide
qu\textquotesingle une VM. On utilise les deux : conteneurs dans une VM
Azure. \\
\end{longtable}

\section{8. Plan
d\textquotesingle apprentissage}\label{8-plan-dapprentissage}

\begin{longtable}[]{@{}
  >{\raggedright\arraybackslash}p{(\linewidth - 4\tabcolsep) * \real{0.1316}}
  >{\raggedright\arraybackslash}p{(\linewidth - 4\tabcolsep) * \real{0.1316}}
  >{\raggedright\arraybackslash}p{(\linewidth - 4\tabcolsep) * \real{0.7368}}@{}}
\toprule\noalign{}
\begin{minipage}[b]{\linewidth}\raggedright
Étape
\end{minipage} & \begin{minipage}[b]{\linewidth}\raggedright
Durée
\end{minipage} & \begin{minipage}[b]{\linewidth}\raggedright
Objectif
\end{minipage} \\
\midrule\noalign{}
\endhead
\bottomrule\noalign{}
\endlastfoot
1 & 1 h & Apprendre les notions de la section 2 (vidéo ou doc k6 « load
testing », bases Docker). \\
2 & 1 h & Lire \texttt{load-test.js} et \texttt{k6-console.txt} ; savoir
commenter chaque ligne du résumé. \\
3 & 2 h & Lire \texttt{test\_reliability.py} scénario par scénario ;
lire les 2 correctifs (fichiers 4 et 5). \\
4 & 1 h & Lancer \texttt{perf} puis \texttt{reliability} sur ton
poste. \\
5 & 2 h & Lire \texttt{docker-compose.yml}, \texttt{pipeline.yml}, le
guide d\textquotesingle installation ; installer le projet sur ton PC en
le suivant. \\
6 & 1 h & Réciter les réponses du guide de soutenance (membre 3) et de
la section 7. \\
\end{longtable}

\section{9. Auto-évaluation}\label{9-auto-uxe9valuation}

\begin{itemize}
\tightlist
\item[$\square$]
  J\textquotesingle explique le p95 et pourquoi on ne regarde pas la
  moyenne.
\item[$\square$]
  Je lis à voix haute un résumé k6 et je dis si le seuil est tenu.
\item[$\square$]
  Je cite les 8 scénarios de fiabilité et les 2 défauts corrigés (cause
  et correctif).
\item[$\square$]
  J\textquotesingle explique ce qui se passe quand l\textquotesingle IA
  est en panne pendant une candidature.
\item[$\square$]
  J\textquotesingle explique les étapes de la CI/CD et
  d\textquotesingle où viennent 20 et 22 jobs.
\item[$\square$]
  J\textquotesingle ai installé le projet sur un autre poste avec le
  guide.
\end{itemize}

\end{document}
